\section{Binary composition and the robust outer circuit}\label{sec:outer-composition}

We now convert the symbol representation into a generalized circuit. Two
features of the conversion require separate arguments. A local Boolean
function may inspect polynomially many bits of a bounded number of symbols;
its implementation must tolerate failures among its own gates. Moreover,
sharing such implementations must not allow one corrupted physical value
to spoil an unbounded number of decisions. We state the first interface,
proved in Section~\ref{sec:inner-gadget}, and then give the construction and
all the counting and synchronization arguments for the second.
The same sequence is used inside a decision gadget in
Section~\ref{sec:inner-gadget}: first obtain replica synchronization without
trusting decisions, then bound unstable comparisons and failed computations,
and finally use the clipped mean displacement to locate an endpoint.
At the outer scale, a completed decision is shared by a group of output
replicas, called a cohort; its internal gates and external read occurrences are charged separately from
each output's own remaining gates.

Throughout this section, a distance between words is relative Hamming
distance, unless a norm is displayed. A constant may depend on previously
fixed accuracies, but never on the End-of-Line input or its length. A
``fixed branch'' of a randomized answer tree means that the earlier
hypothetical answers are fixed when the randomized algorithm is described.
It does not mean conditioning the probability space on receiving those
answers. This distinction applies also to branches never executed.

\subsection{A fixed port and a robust Boolean decision}

The port family will encode every raw symbol in the same binary format.
Its tester lets us recognize symbols close enough to use, and its recovery
radius is fixed before we know how many symbols a corrected local
computation will inspect. Once that arity is fixed, the robust-decision part of
the interface realizes one Boolean output of such a computation. It
provides many designated copies of the answer, so the guarantee remains
useful even when failed gates concentrate on designated outputs.

\begin{theorem}[Standardized ports and robust decisions]\label{thm:port-gadget}
For every fixed positive integer exponent allowance $c_{\rm allow}$, there
are a cutoff $L_{\rm port}$ and a uniform family of injective binary
encodings, for $L\geq L_{\rm port}$,
\[
 \mathcal D:\{0,1\}^{L}\longrightarrow\{0,1\}^{S_0},
 \qquad S_0=\poly(L),
\]
with fixed positive relative distance $d_B$, satisfying \textup{(i)}--
\textup{(ii)} below. The radius $\nu_B$ in \textup{(iii)} is fixed with
this port family; the actual-gadget assertions in that part have the
separate later cutoff specified there.
\begin{enumerate}[label=\textup{(\roman*)}]
\item Every coordinate of $\mathcal D(s)$ is a multilinear polynomial in
the bits of $s$ over $\F_2$, of degree at most $10$. Its coefficient list
is computable in time polynomial in $L$. The degree bound is independent
of $c_{\rm allow}$. Encoding and exact recovery from a fixed positive
relative-error radius are deterministic polynomial-time algorithms.
\item A perfect-completeness tester makes a bounded number of bit queries
and performs bounded Boolean answer processing. Its public randomness has
an enumerable polynomial-size rational support and polynomial bit
complexity. For every fixed $a>0$, there are a fixed $\theta_B(a)>0$
and a cutoff $L_B(a)\geq L_{\rm port}$, allowed to depend on $a$, such
that for every $L\geq L_B(a)$ rejection below $\theta_B(a)$ implies
distance less than $a$ from $\mathcal D(\{0,1\}^L)$.
A word at distance $a$ from an encoding has
rejection at most $C_Ba$, for a fixed $C_B$. Each fixed-branch query has a
bounded-density marginal relative to the uniform distribution on the
block. These assertions include arbitrary words and all hypothetical
branches.
\item There is a constant $0<\nu_B<d_B/4$ determined with this encoding,
before any arity is specified. Subsequently fix an arity bound $k$, a
size factor $C_f$, and a desired output loss $\upsilon>0$. There is a
further cutoff $L_g$, allowed to depend on these fixed choices, such that
all actual-gadget assertions below hold for every
$L\geq\max\{L_{\rm port},L_g\}$. Consider any single-bit Boolean
function of at most $k$ length-$L$ strings, described
by an acyclic fan-in-at-most-two circuit over a fixed Boolean basis with
public constants, of size at most $C_f L^{c_{\rm allow}}$. Some inputs
may be declared public $0^L$.

There is a uniform generalized subcircuit for this function, with a
common count scale $G_0=\poly(L)$ for all padded functions under these
bounds, at most $C_gG_0$ gates, and a set of $W$ designated output wires
produced by its gates, where
\begin{equation}\label{out:designated-count}
 c_gG_0\leq W\leq C_gG_0.
\end{equation}
Here $c_g,C_g>0$ are fixed. Each nonpublic external input is a block of
physical reads $X=3x-1$, with $x\in[0,1]$, compared to a public common
rational threshold $\xi\in[2/5,3/5]$. Every external comparison occurrence
in the complete unfolded subcircuit can be routed separately to an
external replica.
For each external port position, the total number of these occurrences is
at most
\begin{equation}\label{out:port-occurrence}
 C_{\rm occ}G_0/S_0,
\end{equation}
uniformly over $\xi$, the function, and its public parameters. This counts
all hypothetical branches, including unsuccessful and default branches.

There are fixed $\delta_g,\gamma_g,\eps_g>0$ with the following universal
soundness property. Fix logical strings $s_1,\ldots,s_k$ and ideal binary
port words $y_j$ satisfying
$\dist(y_j,\mathcal D(s_j))\leq\nu_B$ for each nonpublic port.
For any subsequently chosen fixed physical margin $\omega>0$, use
\[
 0<\eps\leq\eps_g,\qquad \eps<\omega/1000.
\]
Suppose at most $\gamma_gG_0$ external comparison occurrences fail to
have a physical value on the side of $\xi$ specified by the corresponding
bit of $y_j$, at gap greater than $\omega$. Repeated occurrences need not
have consistent physical values. Then every assignment of the internal
nodes in $[0,1]$ with at most $\delta_gG_0$ failed gates makes all but
$\upsilon W$ designated wires correct near-bits for $f(s_1,\ldots,s_k)$:
at most $\eps$ for zero and at least $1-\eps$ for one.
\end{enumerate}
The port construction and $\nu_B$ are independent of $k,C_f,\upsilon$.
The later constants, the gate-count bounds, and the topology bounds may
depend on these fixed choices, but not on $L,\xi$, or the later margin
$\omega$. The threshold may appear as a rational gate parameter.
Construction and other parameter computations take polynomial time in
$L$ and the supplied description lengths, including the bit length of
$\xi$; all parameter lengths are bounded by a polynomial in these
quantities. In particular an arbitrarily long externally supplied
threshold is charged for its bits.
\end{theorem}

Sections~\ref{sec:inner-algebra} and~\ref{sec:inner-binary} construct the
algebraic and binary ingredients. Proposition~\ref{prop:dummy-port}
proves the fixed encoding, tester, and recovery assertions. The remaining
construction in Section~\ref{sec:inner-gadget} proves the physical occurrence
bound~\eqref{ig:external-load}, gate count~\eqref{ig:gate-count}, and robust
output guarantee~\eqref{ig:extraction-loss}, using the independent geometric
procedures proved in Section~\ref{sec:geometry}. None of these constructions
uses the outer circuit constructed here.

\subsection{Port substitution, validation, and average correction}

Apply Theorem~\ref{thm:outer-representation}, writing its raw encoding as
$R_w=(s_k)_{k\in\Gamma}$, with
\[
 \ell_N=\lceil\log_2(N+2)\rceil,\qquad
 L=(\ell_N+C_{\rm base})^{C_L},\qquad
 |\Gamma|=O(N\ell_N^{O(1)}).
\]
Choose $c_{\rm allow}>c_{\rm raw}+12$, with any additional fixed local
polynomial overhead absorbed in this inequality, and construct
$\mathcal D$ before choosing correction rounds. Put
\begin{equation}\label{out:binary-encoding}
 E_w=(\mathcal D(s_k))_{k\in\Gamma},\qquad m=|\Gamma|S_0.
\end{equation}
The product of the two relative distances is a positive Hamming-distance
bound; write $d_*>0$ for its square root, a lower bound on normalized
Euclidean separation. The distinguished source and its actual successor are publicly computable in polynomial time.

There is a fixed radius $\rho_{\rm dec}>0$ from which these words can be
decoded globally. Indeed, let $r_B>0$ be a port recovery radius and
$r_{\rm raw}>0$ a raw recovery radius. If the whole binary distance is
less than $r_Br_{\rm raw}/2$, at most an $r_{\rm raw}/2$ fraction of port
blocks can be at distance at least $r_B$ from their true block. Decode
all ports, using a fixed default on an invalid or unsuccessful decoder
output, and then use the raw decoder. An incorrect decoding at a bad
site need not be detectable; the raw error bound already charges all
such sites. Fix a smaller radius $\rho_{\rm dec}$ if necessary.

The graph in Theorem~\ref{thm:outer-representation} has indegree and
outdegree at most one, no self-edges, and no isolated vertices. Write
$w_*$ for its distinguished source, which has an actual successor, and
call every other missing-neighbor vertex a desired endpoint. The
geometric construction applies to this encoded graph. For an edge
$u\to v$, its four-block polygon successively changes the second, third,
first, and third blocks:
\[
 \begin{aligned}
 (E_u,E_u,0,0)&\longrightarrow(E_u,E_v,0,0)
 \longrightarrow(E_u,E_v,1,0)\\
 &\longrightarrow(E_v,E_v,1,0)
 \longrightarrow(E_v,E_v,0,0).
 \end{aligned}
\]
Here a scalar in a block denotes its constant repetition. The corners
are rounded, and a segment in the fourth block enters the distinguished
source from $(E_{w_*},E_{w_*},0,2)$.
Theorem~\ref{thm:geometry} constructs a displacement field on
\[
 \mathcal B=[-1,2]^{4m},\qquad
 \norm{U}_{\rm b}^2=\frac1m\sum_i U_i^2,
\]
where all distances to subsets of $\mathcal B$ use this norm. Each point
of its endpoint set $\mathcal E$ has first block exactly $E_w$ for a desired
endpoint $w$; the top of the extra segment is excluded. For any sufficiently
small fixed rational $\tau>0$, chosen to meet the validation and
global-recovery bounds in \eqref{out:tau}, there are constants
$B_1,c_*,\beta_0>0$ and a field
$v$ such that, after increasing $B_1$ to an integer,
\begin{equation}\label{out:geometry-displacement}
 \norm{v(U)}_\infty\leq B_1-1,\qquad
 \norm{\clip_{\mathcal B}(U+\beta v(U))-U}_{\rm b}\geq c_*\beta
\end{equation}
whenever $\dist(U,\mathcal E)>10\tau$ and $0<\beta\leq\beta_0$.
These constants precede estimator accuracy and the step $\beta$.
The circuit will keep many replicas of each physical coordinate and
force their mean vector close to $\mathcal E$; rounding the first block
will then permit global recovery.

To approximate the field locally, the numerical procedure of
Lemma~\ref{lem:geometric-procedures} rounds the first two blocks at one
common threshold $\xi\in[2/5,3/5]$. Each rounded word is a seed for
validation. On good decisions an accepted seed must identify one fixed
reference vertex, independent of the random acceptance decision and of
the requested confidence or numerical accuracy; a seed in the intermediate
region may be accepted or rejected. If one of the first two physical
blocks is within $5\tau$ of a valid word, its seed must be accepted with that
reference on good decisions for every allowed threshold. The procedure
then asks for incident-edge flags, bits of each reference and its actual neighbors,
and bits of the distinguished source and its successor. Each word-bit
request uses either the coordinate matching the field component being
estimated or a preallocated uniform coordinate for a bounded scalar
statistic. Each reference and its neighbors determine a bounded list of
incident path pieces; a fixed mesh and estimates of their scalar data
approximate the field near those pieces. Scalar reads use fixed-width
quantization with a public offset. We next supply these calls. Their
number is fixed once the geometric mesh, scalar sample counts and
numerical tolerances have been chosen, before the smaller bit-error
budgets are assigned.

\begin{lemma}[Protecting functional inputs]\label{out:flags}
There are constants $0<a_1<a_2<\nu_B$, fixed before the arity of any
functional invocation, with the following property. A bounded Boolean
flag routine on an arbitrary binary port block can have arbitrarily
small fixed failure probability, unflagging every block within $a_1$ of
an encoding and flagging every block at distance at least $a_2$ from all
encodings. In a functional invocation replace flagged ports by public
$0^L$. Except on the union of the completed flag routines' failure
events, every nonpublic input of every selected invocation is within
$\nu_B$ of its intended encoding.
\end{lemma}
\begin{proof}
Choose $a_2<\nu_B$ and then $a_1$ so small that
$C_Ba_1<\theta_B(a_2)/4$. Estimate port-tester rejection by a fixed
number of independent trials and compare to $\theta_B(a_2)/2$.
Bounded independent variables give any fixed confidence by increasing
the fixed sample count. In the interval between the two promises either
Boolean result is permitted. An unflagged block on the good event is
within $a_2$ of an encoding, unique because $a_2<d_B/4$; take its logical
string as the intended input. A flagged input is explicitly a public
zero, so it requires no external near-port premise. A union bound is
applied to the finite list of completed flag promises, after that list
has been fixed. Their internal tester answer trees are not counted as
new logical promises.
\end{proof}

Fix a physical vector $U$ and initially fix its rounding threshold and
scalar offset. A rounded seed $y=(y_k)_{k\in\Gamma}$ is now a fixed
binary word; use bit one at equality with the threshold. For analysis, define a global nearest-symbol convention
$s^*(k)$ at every site within $a_2$ of a port encoding, and choose a
fixed arbitrary string elsewhere. An ideal functional call computes
its Boolean function on these symbols at unflagged inputs and on
$0^L$ at flagged inputs. Within one correction cone, deduplicate repeated
raw sites and share their flag and initial logical value. Different
cones may make different choices at uncertain sites.

\begin{lemma}[Binary validation]\label{out:validation}
There are fixed $0<\rho_{\rm comp}<\rho_{\rm acc}$, with
$\rho_{\rm acc}$ smaller than the uniqueness radius of
\eqref{out:binary-encoding}, and a finite ideal Boolean validation
procedure such that, with arbitrarily high fixed probability, acceptance
implies distance less than $\rho_{\rm acc}$ from a unique valid $E_w$,
and distance less than $\rho_{\rm comp}$ implies acceptance with the
true reference. The procedure uses port flags, port tests, and bounded
raw-checker functional calls. It remains valid with averaged stationary
expander-walk sampling as described below.
\end{lemma}
\begin{proof}
Choose $\rho_{\rm acc}$ small enough that
$\rho_{\rm acc}/a_1$ is within the raw uncertainty-correction radius.
Choose $\alpha<\rho_{\rm acc}/4$ and take the raw gap
$0<\theta_\alpha\leq1$ from Theorem~\ref{thm:outer-representation}.
Writing $q=|\Pi|$ for its bounded checker-neighborhood size, choose
$0<a<a_1$ and $b_0>0$ satisfying
\begin{equation}\label{out:validation-constants}
 \alpha+a+b_0<\rho_{\rm acc},\qquad
 qb_0<\theta_\alpha/16,
 \qquad p_{\max}=b_0\theta_B(a)>0.
\end{equation}
Decrease $p_{\max}$ if necessary. Let $p$ be the average port-tester
rejection over all sites. Estimate $p$, and independently simulate the
raw checker at sampled centers using the flagged/substituted inputs.
Accept when the first estimate is at most $p_{\max}/2$ and the simulated
raw rejection estimate is at most $\theta_\alpha/2$.

On statistics of error less than $p_{\max}/4$, acceptance gives
$p<p_{\max}$. Consequently the fraction of sites not within $a$ of any
port encoding is less than $b_0$. At a center whose entire neighborhood
avoids these sites, good flag promises leave all inputs unflagged and
make the simulated checker exactly the checker on the global word
$s^*$. Exceptional centers occupy at most $qb_0$. In the analysis,
also estimate the true checker result on $s^*$ to accuracy
$\theta_\alpha/8$ and the exceptional-center fraction to accuracy
$\theta_\alpha/16$. These are fixed bounded statistics; the algorithm
need not compute the hypothetical true statistic. On their good events
and the completed flags' good events,
\[
 \Pr[\text{raw checker rejects }s^*]
 \leq\theta_\alpha/2+\theta_\alpha/8+qb_0+\theta_\alpha/16
 <\theta_\alpha.
\]
Raw soundness makes $s^*$ $\alpha$-close to a valid $R_w$. Charging all
raw-disagreement sites and all non-$a$-close sites at cost one yields
\[
 \dist(y,E_w)<\alpha+a+b_0<\rho_{\rm acc}.
\]
The uniqueness radius makes this reference independent of subsequent
random acceptance decisions for the fixed $y$.

Conversely, if $\dist(y,E_w)<\rho_{\rm comp}$, then
$p\leq C_B\rho_{\rm comp}$. By Markov's inequality at most
$\rho_{\rm comp}/a$ of the ports are at distance at least $a$ from their
true block. Outside their $\Pi$-neighborhoods the simulated checker is
the perfectly accepting raw checker. Choose $\rho_{\rm comp}$ so small
that both means are strictly below one quarter of their acceptance
thresholds; choose the sampling errors with the remaining slack. Flag
error probabilities are set after the sample counts. This proves both
assertions. The proof requires concentration of the sampled exceptional
fraction when $p$ is small. It never requires an arbitrarily long batch
to avoid every exceptional site.
\end{proof}

Choose the geometric radius so that
\begin{equation}\label{out:tau}
 4(10\tau)^2<\rho_{\rm dec},\qquad
 \frac{(5\tau)^2}{(2/5)^2}<\rho_{\rm comp}.
\end{equation}
The second inequality supplies physical validation completeness:
every incorrectly rounded coordinate of a binary word costs at least
$(2/5)^2$ in squared distance for every allowed threshold. Fix the
resulting $B_1,c_*,\beta_0$ and fix $\kappa=1/100$, the replica gap used
below. Set the unscaled estimator error target with slack, smaller than
$c_*\kappa/100$ in normalized root-mean-square norm.

\begin{lemma}[Average correction of every relevant word role]
\label{out:average-correction}
After fixing the high-level mesh and statistical call counts, boundedly
many raw correction rounds suffice to supply every required incident
flag and every requested encoding bit, with any prescribed fixed average
error. The average is over a uniform output component or its
uniform-marginal statistical requests. The conclusion holds
unconditionally for every true reference that can be accepted on good
validation decisions, before selecting pieces or numerical winners.
All resulting single-bit functions fit the fixed exponent allowance of
Theorem~\ref{thm:port-gadget}.
\end{lemma}
\begin{proof}
Fix a rounded seed that can be accepted on good validation, and its
unique reference $R_w=(s_k)$. Put
\[
 U_0=\{k:\dist(y_k,\mathcal D(s_k))>a_1\},
 \qquad |U_0|/|\Gamma|\leq\rho_{\rm acc}/a_1.
\]
Outside $U_0$, good flags unflag the port and its unique nearest logical
string is $s_k$. Within $U_0$, substitution and intermediate flag choices
may supply arbitrary values. In each deduplicated cone extend those
initial values to a whole raw word agreeing with $R_w$ outside $U_0$.
The raw uncertainty theorem gives a residual set $U_t$, depending only
on $U_0$, with arbitrarily small fixed density after sufficiently many
fixed rounds $t$. Its guarantee holds simultaneously for every such
extension. Thus different cones and correlated intermediate flags cause
no additional exceptional support. A corrected symbol is right off
$U_t$; a local proposal and edge flag are right off a fixed permutation
expansion of $U_t$.

Compute an edge flag at a uniform-marginal center. For a requested own
or neighboring encoding bit, compute the appropriate raw output symbol
and then the requested coordinate polynomial of $\mathcal D$. Matching
coordinates are uniform in the output average; the statistical requests
have the same center marginal. Each fixed permutation expansion
therefore costs at most its size times $|U_t|/|\Gamma|$. Set $t$ for
the already fixed list of high-level true roles and calls. No union over
all possible intermediate flag patterns is necessary.

The raw simulation uses at most $C_tL^{c_{\rm raw}}$ Boolean gates, with
exponent independent of $t$. A degree-at-most-ten Boolean polynomial has
at most $\sum_{j=0}^{10}\binom Lj=O((L+1)^{10})$ monomials. Compute each
monomial by a bounded-fanin AND tree and their sum by XOR gates expanded
in the fixed basis. Coefficients and the requested bit's description are
publicly generated. The extra bit-index and padding work has a fixed
polynomial exponent. Our initial choice of $c_{\rm allow}$ leaves
slack for this work. The number of raw ports and the multiplicative
factor may increase with $t$, but are now fixed. The port encoding and
its radius do not change.
\end{proof}

The completed port flags are now amplified for the finite list of
logical promises in validation and these correction cones. Conditional
on a sampled site, they use fresh tester randomness. Make the probability
that any selected invocation has a nonpublic input outside $\nu_B$
arbitrarily small. This protection concerns arbitrary rounded seeds,
including seeds undergoing validation, and does not require exclusion
of the geometric endpoint neighborhood. Uncertain flag results only
choose which already uncertain inputs are substituted.

\subsection{Small-support ideal tapes and complete branch counts}

\begin{lemma}[Averaged expander sampling]\label{out:walk}
The statistics needed above can be sampled with polynomial-in-$L$
random support per output coordinate, with any desired fixed confidence
in the average over output coordinates.
\end{lemma}
\begin{proof}
Write an output coordinate as $(b,k_0,z_0)$ with $b\in[4]$,
$k_0\in\Gamma$, and $z_0\in[S_0]$. Use fixed-length labeled walks in the
outer Cayley expander starting at $k_0$. Fresh bounded segments and
auxiliary uniform positions in $[S_0]$ supply all center and word
statistics. Averaging over uniform $k_0$ makes every walk position
stationary. If the mean-zero operator norm is at most $\lambda<1$,
bounded centered functions at times $s,t$ have covariance at most
$C\lambda^{|s-t|}$. Independent auxiliary randomness only changes the
bounded diagonal terms and replaces off-diagonal functions by their
conditional means. Therefore a batch of $q_s$ samples satisfies
\[
 \operatorname{Var}\left(\frac1{q_s}\sum_{j=1}^{q_s} Z_j\right)
 \leq\frac{C}{q_s^2}\sum_{s,t=1}^{q_s}\lambda^{|s-t|}
 \leq\frac{C(1+\lambda)}{(1-\lambda)q_s}.
\]
Chebyshev's inequality and a finite union give any fixed confidence by
choosing a sufficiently large fixed $q_s$. The argument permits distinct
bounded functions at the different slots and separate batches for
physical blocks.

Apply it to port rejection, the fixed global raw-checker and exceptional
indicators in validation, and each actual reference candidate's
statistics in the geometric procedure. At fixed rounded seeds those
reference vertices are fixed whenever good acceptance is possible.
Analyze all their roles before a numerical minimum or uncertain
acceptance chooses which ones to use. The concentration is not
conditioned on the realized answer path. Each walk has a fixed number
of fixed-degree label choices. Each remaining choice has polynomial-in-
$L$ support, and only a fixed number of choices is used. Their joint
support per output coordinate is consequently polynomial in $L$.
\end{proof}

The preceding routines now meet the local-access contract of
Definition~\ref{geo:access-contract}. Validation supplies fixed unique
references, average correction supplies the matching and statistical bit
requests, and the walk lemma supplies concentration for each underlying
true candidate. The geometric procedure reserves at most eight true word
roles and thirty-two piece roles, as enumerated in
Section~\ref{geo:candidates}. Missing roles are omitted; duplicate roles
may be retained.
For fixed $U$ outside the endpoint neighborhood and each fixed admissible
binary and scalar-quantization offset, it gives arbitrarily small
prescribed $4\E_{i,\mathrm{trial}}|\widehat v_i-v_i(U)|^2$, with $i$
uniform among the $4m$ components. All bit and flag errors and all
candidate statistics are analyzed before a numerical winner is selected.
On the resulting good event, Lemma~\ref{geo:stability} controls every
sufficiently near-minimizing candidate in each component before selection, so a winner correlated with
the requested component is harmless. The needed bit-error bounds remain
average bounds over the allocated requests.

Apart from the supplied calls, the procedure uses bounded comparison and
Boolean-answer processing, public rational numerical tables, and a
signed dyadic output in $[-B_1,B_1]$, including on every bad branch.
A fixed-width scalar-quantization grid with an independently randomized
offset suffices. Refining the finite grids of offsets changes comparison
thresholds, not the number of bins, answers, or slots. All these
statements concern the unscaled field.

This outer application uses the accuracy and finite-role clauses
\textup{(i)}--\textup{(iii)} of
Lemma~\ref{lem:geometric-procedures} with completed ideal Boolean
functional calls. Its optional physical-read clause \textup{(iv)} is
not applied to the full outer gadget expansion. The fixed count
$H_{\rm out}$ below counts direct outer read slots, and $B_{\rm out}$
counts ideal functional invocations. Theorem~\ref{thm:port-gadget} and
Lemma~\ref{out:routing} separately bound every physical read created by
substituting those invocations.

Applying parts \textup{(i)}--\textup{(iii)} of
Lemma~\ref{lem:geometric-procedures} at this ideal-oracle level, with
the preceding accuracy choices, we obtain ideal dyadic outputs satisfying,
for every eligible fixed $U$ and every fixed admissible offset,
\begin{equation}\label{out:ideal-accuracy}
 4\E_{i,\mathrm{tape}}\abs{\widehat v_i-v_i(U)}^2
       <(c_*\kappa/100)^2.
\end{equation}
All outputs remain in $[-B_1,B_1]$ on bad trials. The good-port failure
probability can independently be made arbitrarily small on this
high-level call list. It is useful to record the order explicitly:
\begin{enumerate}[label=\textup{(\arabic*)}]
\item Fix geometric roles, mesh, scalar counts and numerical tolerances.
\item Choose validation confidence and raw correction rounds for this
finite logical-call list; fix the number of ports in every cone.
\item Amplify each completed flag for its prescribed failure probability.
Amplification adds tester reads, not new logical promises.
\item Only now unfold the complete outer answer tree, leaving functional
calls unexpanded and obtaining a fixed number $H_{\rm out}$ of direct
read slots and $B_{\rm out}$ of hypothetical functional invocations.
\item Choose $\upsilon$ for these $B_{\rm out}$ invocations and apply
Theorem~\ref{thm:port-gadget} with the now fixed arity and size factor.
\end{enumerate}
This order does not define a bit accuracy by a union bound over the very
answer tree created by amplifying that accuracy.

At the ideal-oracle level, unfold every flag pattern, substitution
pattern, quantized bin and hypothetical answer. Functional-call nodes
remain unexpanded: each contributes one completed Boolean oracle answer,
and its physical implementation is separately accounted for below.
Thus $H_{\rm out}$ counts only direct read slots and $B_{\rm out}$ counts
functional invocations. Supply randomness to unused nodes too. Earlier
answers that determine an invocation's function or addresses are fixed
publicly in that copy; they are not additional fragile input wires.
Thus all port functions remain Boolean functions on their declared raw
ports. A wrong numerical answer selects another reserved role or a
finite table case. It cannot become an arbitrary hidden word address.
Each hypothetical raw site is $\pi(k_0)$ for $\pi$ in a fixed finite set
$\mathcal P$: compose walk-label moves, raw-neighborhood permutations
and correction-cone moves. This set is independent of later offset-grid
refinement. Direct scalar and bit slots use matching $z_0$, independently
sampled uniform within-block positions, or the port tester's smooth
queries. For fixed $k_0$ and fixed offsets, summing their unconditioned
position densities over the finite complete tree gives a constant.

The tape law in the next lemma includes both late offset grids---the
binary threshold and the independent scalar-quantization offset---together
with all internal choices at the hypothetical nodes. The grids are fixed
before this final law is rationalized. The preceding accuracy estimates
hold for each fixed offset, and hence for this joint law; the finite grids
multiply its support bound only by a fixed factor.

\begin{lemma}[Rational tapes and their read loads]\label{out:tapes}
The true tape law can be replaced deterministically, at each output
coordinate, by a list of a common square length $T_{\rm tap}=\poly(L)$,
with total variation at most $\eta_0/S_0$, for any prescribed fixed
$\eta_0>0$. The lists use only choices in the original support. They
preserve \eqref{out:ideal-accuracy} and any bounded rare-event estimate
with arbitrarily small fixed slack. Across these lists, every base
physical coordinate is read by at most $C T_{\rm tap}$ direct
comparison slots in the complete unfolded outer trees.
\end{lemma}
\begin{proof}
Let $Q(L)$ bound the enumerable rational support of every tape law.
For a common square $T\geq Q(L)S_0/\eta_0$, round each mass $Tp_j$
down to an integer and distribute the remaining copies among nonzero
support entries. Each mass error is at most $1/T$, and total variation
is at most $Q(L)/T\leq\eta_0/S_0$. Floors, remainders and the ordering
of entries are computed with exact rational arithmetic in polynomial
time. The probability descriptions have polynomial bit complexity.
A function in $[0,A]$ changes expectation by at most $A$ times total
variation, proving the bounded-statistic assertions.

For the load assertion, explicitly list all potential direct slots
$\mathcal A$, including their fixed earlier-answer paths; set
$H=|\mathcal A|$. If a slot is omitted or supplied a public value, its
read subdistribution is dominated by that of its potential slot; it is
never renormalized on being selected. Fix a target $(b',k',z')$ and a
possible $\pi\in\mathcal P$. Its preimage $k_0=\pi^{-1}(k')$ is fixed.
For every fixed offset, the true law gives
\begin{equation}\label{out:direct-density}
 \sum_{z_0=1}^{S_0}
 \E[\text{number of direct reads at }(b',k',z')]
 \leq C,
\end{equation}
where $C$ is the sum of the unconditioned fixed-slot density constants.
Matching coordinates contribute one, uniform sampled coordinates
contribute their bounded counts, and tester queries contribute their
smoothness constants. Branches and seed roles add only a fixed factor.
A position-count random variable lies in $[0,H]$. Changing the law for
one output by total variation $\eta_0/S_0$ adds at most
$H\eta_0/S_0$ to its expectation; summing over $z_0$ adds only
$H\eta_0$. Sum the fixed preimages and physical roles and multiply by
$T_{\rm tap}$. This is the asserted $O(T_{\rm tap})$ load.
A constant total-variation tolerance without the $S_0$ denominator
would not give this bound.
\end{proof}

The occurrence bound \eqref{out:port-occurrence} for each substituted
gadget uses the same full-tree principle at its own scale. Its tape
approximation must control loads both on its internal word, of length
$m_{\rm in}$, and on an external port, of length $S_0$. Section~\ref{sec:inner-gadget}
therefore uses total variation $\eta/\max\{4m_{\rm in},S_0\}$; the
calculations \eqref{ig:external-load}--\eqref{ig:internal-load} give
$O(G_0/S_0)$ external occurrences per port position and $O(R_{\rm in})$
internal reads per base coordinate, where $R_{\rm in}$ is the gadget's
replica count and $G_0=4m_{\rm in}R_{\rm in}$.
This accounts for unequal polynomial block lengths without conditioning
on an executed branch.

\subsection{Cohorts, individual read routing, and gate realization}

At each base component $i\in[4m]$, allocate one cohort for every tape in
its list. A cohort has $K_c$ output replicas, where $K_c$ is a square
satisfying
\begin{equation}\label{out:cohort-size}
 \max\{1,C_gG_0\}\leq K_c\leq C_KG_0
\end{equation}
for a fixed $C_K$; increasing $C_g$ if necessary also covers all
designated-wire counts. The total number of replicas per component is
$R=T_{\rm tap}K_c$, again a square. Each cohort contains its complete
bounded list of functional invocations, sharing these invocations
among its $K_c$ outputs. For each invocation, assign its $W$ designated
wires to these outputs by round robin. Since $W\leq K_c$,
\begin{equation}\label{out:round-robin}
 \#\{\text{uses of any designated wire}\}
 \leq\lceil K_c/W\rceil\leq 2K_c/W.
\end{equation}
A fraction $\upsilon$ of wrong designated wires therefore affects at
most $2\upsilon$ of the cohort outputs. The remaining gates, called a
direct cone, are separately allocated for each output replica.

\begin{lemma}[Bounded physical read multiplicity]\label{out:routing}
Every physical replica can be made the input of at most a fixed number
of comparison occurrences, counting all outer and inner hypothetical
branches. The construction is deterministic and uses no shared rounded
bit between different comparison occurrences.
\end{lemma}
\begin{proof}
Fix a physical block and a raw site $k'$. For each of the fixed possible
address permutations, a functional port requesting $k'$ has a fixed
$k_0$ preimage. There are $S_0$ output bit positions and $T_{\rm tap}$
tapes at that preimage, and only a bounded number of invocation/port
roles. Thus at most $C S_0T_{\rm tap}$ invocation ports request $k'$.
By \eqref{out:port-occurrence}, any position in that site is used at most
\begin{equation}\label{out:gadget-load}
 C S_0T_{\rm tap}\,\frac{C_{\rm occ}G_0}{S_0}
       \leq C' T_{\rm tap}K_c=C'R
\end{equation}
times. Lemma~\ref{out:tapes}, followed by replication over the cohort,
gives $O(T_{\rm tap}K_c)=O(R)$ direct comparison occurrences at each
base coordinate. Enumerate all these occurrences and assign successive
ones round robin to its $R$ physical replicas. Each actual replica has
bounded resulting multiplicity. Each comparison is performed in its
own gate cone directly on that replica's $x$ value. In particular, an
error in one comparison wire is not broadcast to other inner cones.
\end{proof}

Write the physical replica values and their means as
\[
 X_{ir}=3x_{ir}-1,\qquad \bar X_i=R^{-1}\sum_rX_{ir},\qquad
 \norm{Y}_{\rm rep}^2=\frac1{mR}\sum_{i=1}^{4m}\sum_{r=1}^R Y_{ir}^2.
\]
A vector constant over replicas has replica norm equal to its base norm.
Let $J$ be orthogonal projection onto these constant-over-replica
vectors, so $JX=\bar X$ with the usual identification.

For each $i$, average the replica labels using the half-lazy
Margulis--Gabber--Galil operator $A_{\rm rep}$ on
$(\mathbb Z/q\mathbb Z)^2$, $q=\sqrt R$. Its sixteen equally weighted
labels consist of eight identity labels and the eight moves
\[
 (a\pm2b,b),\quad(a\pm(2b+1),b),\quad
 (a,b\pm2a),\quad(a,b\pm(2a+1)),
\]
with multiplicities retained. For every positive integer modulus the
normalized eight-move operator is symmetric with nonconstant numerical
radius at most $5\sqrt2/8$; see the primary formal proof
\cite[\texttt{mgg\_numerical\_radius}]{MGGAFP}. Thus
\begin{equation}\label{out:replica-gap}
 A_{\rm rep}J=J,\qquad
 \norm{A_{\rm rep}(I-J)Y}_{\rm rep}
       \leq(1-\kappa)\norm{(I-J)Y}_{\rm rep},
 \qquad \kappa=1/100.
\end{equation}
The modulus need not be prime. Loops and parallel labels are kept as
individual averaging inputs. The mean-zero space for modulus one is
zero. Every averaging connection has bounded multiplicity as well.

We specify the direct cones using exactly the gate types in
Definition~\ref{def:gcircuit}. This also specifies how public constants
and finite numerical processing are charged. Every needed constant
node is the fresh output of a Constant gate in the cone using it, unless
its value is incorporated as a legal Scale parameter or into a Boolean
truth table. A physical comparison at threshold $t\in[-1,2]$ is a
Constant gate with parameter $(t+1)/3$ and a Less gate on the physical
replica's $x$ node and that constant, with orientation chosen for the
required bit. Outside the physical range its answer is a public
constant. On a satisfied comparison, an ideal mean gap greater than
$4\omega$, an actual read deviation at most $\omega$, and
$\eps<\omega/1000$ leave a strict Less antecedent even after the
constant error: the node-scale gap is at least $\omega-\eps>\eps$.

Boolean finite tables, flag processing, and quantization-bin processing
use And, Or and Not with fan-in at most two. If their inputs are correct
near-bits, each gate produces the same correct near-bit interval,
without accumulating Boolean error. In particular the weak inequalities
at $\eps$ and $1-\eps$ are sufficient. A selection among hypothetical
invocations is implemented as a Boolean multiplexer: And each candidate
answer with its selector and Or the resulting terms. An inactive
selector is a low near-bit and forces that term low even if the unused
candidate answer is arbitrary in $[0,1]$. Thus unsuccessful unused
branches are shielded, without assuming that their outputs are bits.
All finite answer tables have fixed size, although their public entries
may depend on the tape and threshold.

Represent the ideal dyadic field value by
\begin{equation}\label{out:signed-value}
 \widehat v=B_1(p^*-n^*),\qquad
 p^*,n^*\in[0,1],\quad p^*n^*=0,
\end{equation}
using a fixed number of binary digits; enlarge the dyadic precision
within the earlier estimator slack. Scale the near-bit digits by their
dyadic weights and combine them by clipped Add gates to obtain $p,n$.
Compute the $x$-average by sixteen Scale gates of weight $1/16$ and a
bounded Add chain. All exact partial sums are in $[0,1]$. Choose a
positive rational step with
$t_\beta=\beta B_1/3\leq1$. Scale $p,n$ by $t_\beta$, add the positive
term to the computed average, then subtract the negative term. The last
output, or a final Copy output, is $x_{ir}$ itself. When the ideal signs
are exclusive this expression is exactly
\[
 \frac{1+\clip_{[-1,2]}((A_{\rm rep}X)_{ir}
                         +\beta\widehat v)}3.
\]
Indeed if the negative term vanishes the only possible clipping is at
the upper end; if the positive term vanishes it is at the lower end.
The affine change of coordinates commutes with the stochastic average.

\begin{lemma}[Coarse and fine direct-cone bounds]\label{out:direct-cone}
There is a fixed $C_{\rm ar}$, determined by the finite direct-cone
shape before $\beta$ is filled in, such that every entirely
$\eps$-satisfied direct cone obeys
\begin{equation}\label{out:coarse-pointwise}
 |X_{ir}-(A_{\rm rep}X)_{ir}|
       \leq2\beta B_1+C_{\rm ar}\eps.
\end{equation}
No correctness of its Boolean or functional inputs is required.
If its direct comparisons and selected invocation answers are correct
near-bits for the ideal run on $\bar X$, then also
\begin{equation}\label{out:fine-pointwise}
 |X_{ir}-\clip_{[-1,2]}((A_{\rm rep}X)_{ir}
                +\beta\widehat v_{ir}(\bar X))|
       \leq C_{\rm ar}\eps.
\end{equation}
\end{lemma}
\begin{proof}
All hidden nodes, including $p,n$, remain in $[0,1]$ even when their
logical values are wrong. Satisfied final Scale gates therefore make
each signed term at most $t_\beta+\eps$. The absolute change from the
computed average is at most their sum plus the final arithmetic errors.
The averaging chain is within $C\eps$ of its exact average. Returning
to physical units gives \eqref{out:coarse-pointwise}, after enlarging
$C_{\rm ar}$. For the fine bound, induction through the shielded
Boolean selections supplies the correct digits in
\eqref{out:signed-value}. Scale is nonexpansive and clipped Add and
Subtract are Lipschitz with respect to the sum of their input errors.
For an explicit bound, let $J_d$ be the number of dyadic digits in
one magnitude. A magnitude uses $J_d$ Scale and $J_d-1$ Add gates;
its error is at most $(3J_d-1)\eps$. The sixteen Scale and fifteen Add
gates for the average contribute at most $31\eps$. The two final
Scale gates and the final Add/Subtract pair therefore contribute a
total error at most $(6J_d+33)\eps$ in node units, or
$(18J_d+99)\eps$ in physical units. An optional final Copy adds
$3\eps$. The coarse bound's arithmetic error is at most $108\eps$
including that Copy. Thus $C_{\rm ar}=18J_d+108$ suffices for this
implementation, uniformly in all Scale parameters in $[0,1]$.
The exclusive-sign identity gives \eqref{out:fine-pointwise}.
\end{proof}

Every invocation owns a disjoint gate set, and every direct cone owns a
disjoint gate set. Internal outputs are fresh; designated and external
wires are read-only fan-outs. The cyclic physical replica output is
produced by the final gate of its own unique cone and by no other gate.
There is no promised free constant node. Let
\begin{equation}\label{out:gate-count}
 K_{\rm out}=4mR,
 \qquad |T|\leq C_{\rm tot}K_{\rm out}.
\end{equation}
The second assertion follows because a cohort has a bounded number of
invocations, each with $O(G_0)=O(K_c)$ gates, and $K_c$ direct cones of
bounded size. Each designated output is included in its invocation's
gate count. In particular, concentrating failures on extraction gates
receives the same charge as any other failures.

\subsection{All failure and spoilage charges}

We next choose the late offsets. Use a uniform rational grid in
$[2/5,3/5]$ for $\xi$ and, independently, a uniform rational grid modulo
the fixed scalar-bin width for the quantization offset. For every fixed
physical value, the probability of lying within $4\omega$ of a binary
threshold or any relevant shifted scalar boundary is at most
\begin{equation}\label{out:anticoncentration}
 C_{\rm bin}(D_{\rm grid}^{-1}+\omega),
\end{equation}
where $D_{\rm grid}$ is a lower bound on the two grid sizes and
$C_{\rm bin}$ depends only on the already fixed widths and bin count.
This follows by counting grid points in intervals of total length
$O(\omega)$, adding at most a fixed number of boundary grid points.
The number of comparison slots is unchanged when the grids are refined.
The ideal accuracy estimate is conditional on every allowed offset, so
refinement cannot invalidate it.

Call a routed comparison occurrence \emph{spoiled} if its mean value is
within $4\omega$ of its prescribed boundary or if its physical replica
value differs from its mean by more than $\omega$. Otherwise its input
has the prescribed ideal side with physical gap greater than $\omega$.
For direct cones, spoiled comparisons are charged to their output. For
functional invocations, count the occurrences themselves, including
all hypothetical branches.

\begin{lemma}[Deterministic charges for arbitrary assignments]
\label{out:charges}
There are fixed counting constants with the following properties.
For an assignment with at most $\delta|T|$ failed gates:
\begin{enumerate}[label=\textup{(\roman*)}]
\item At most a $C_{\rm tot}\delta$ fraction of direct cones contains a
failed gate.
\item Cohorts containing an invocation with more than $\delta_gG_0$
failed gates account for at most $C\delta/\delta_g$ of all outputs.
\item The total number of excessive-deviation comparison occurrences,
divided by $K_{\rm out}$, is at most
$C\norm{X-\bar X}_{\rm rep}^2/\omega^2$.
\item The averaged number of ambiguous gadget comparison occurrences,
divided by $K_{\rm out}$, is at most
$C(D_{\rm grid}^{-1}+\omega+\eta_0)$.
The fraction of direct cones with an ambiguous comparison has the same
form of bound. Here the averaging over tapes is the actual empirical
frequency of the circuit, and the bounds hold for every fixed mean
vector.
\item Cohorts containing an invocation with more than $\gamma_gG_0$
spoiled external occurrences account for at most $C/\gamma_g$ times
the total normalized gadget-spoil count.
\end{enumerate}
\end{lemma}
\begin{proof}
Disjointness of direct gate sets proves (i). The number of invocations
violating the threshold in (ii) is at most
$\delta|T|/(\delta_gG_0)$. Charging $K_c$ outputs for each and using
\eqref{out:cohort-size} and \eqref{out:gate-count} proves (ii), even if
the same cohort is charged repeatedly. For (iii), Markov's inequality
gives at most
$mR\norm{X-\bar X}_{\rm rep}^2/\omega^2$ excessive physical replicas.
Lemma~\ref{out:routing} bounds their total comparison multiplicity by
a constant. The normalization $K_{\rm out}=4mR$ only improves the bound.

For (iv), fix a port of an invocation at site $k$ and threshold $\xi$.
Equation~\eqref{out:port-occurrence} bounds its ambiguous occurrences by
\[
 C_{\rm occ}G_0\,
 \frac{\#\{z:|\bar X_{kz}-\xi|\leq4\omega\}}{S_0}.
\]
Every site is an image of a uniform $k_0$ under a member of the fixed
permutation list. Sum over that list and all invocation ports under the
unconditioned tape law, before selecting any answer branch. Averaging the offset and using
\eqref{out:anticoncentration} bounds the expected number per trial by
$C G_0(D_{\rm grid}^{-1}+\omega)$. The full occurrence bound is uniform
in threshold-dependent function data. Tape rationalization changes a
count bounded by $CG_0$ by at most $CG_0\eta_0/S_0$. Dividing each
trial's charge by its $K_c$ output replicas gives (iv). Direct slots use
the full-tree density argument in \eqref{out:direct-density}, and the
same grid count; their count per trial is bounded by $H_{\rm out}$.
This reasoning applies to every fixed eventual mean, rather than a
mean selected before construction. Finally, (v) is the same threshold
counting proof as (ii), with spoiled occurrences replacing failed gates.
\end{proof}

\subsection{Parameter order and synchronization before decisions}

We give an acyclic choice of all remaining constants. Reserve enough
slack in \eqref{out:ideal-accuracy} that after tape rationalization the
ideal error norm is at most $c_*\kappa/20$. Choose $q_0>0$ so small that
\begin{equation}\label{out:q0}
 20(B_1+1)\sqrt{q_0}<c_*\kappa/20.
\end{equation}
The geometric and statistical accuracies, correction rounds, and
completed flags are chosen as above to put the ideal good-port failure
probability below, for example, $q_0/20$. These are logical-call budgets.
After the complete outer tree is fixed, choose
$2B_{\rm out}\upsilon<q_0/20$ and obtain the gadget constants. Increasing
$K_c$ for these gadgets weights every trial equally and does not
magnify its ideal failure probability.

Now choose fine fixed offset grids, then a small positive rational
$\omega$, then a small tape tolerance $\eta_0$, and then a desired
synchronization radius $d_{\rm sync}>0$. By
Lemmas~\ref{out:tapes} and~\ref{out:charges}, they can be chosen so that,
whenever $\norm{X-\bar X}_{\rm rep}\leq d_{\rm sync}$, the sum of the
following empirical output fractions is less than $q_0/2$: failure of
the ideal good-port event, direct ambiguous/excessive comparisons,
and cohorts having an invocation with too many spoiled occurrences.
For example assign $q_0/20$ to each grouped term, with the
$C/\gamma_g$ factors included before choosing the responsible parameter.
Ideal squared error is bounded, so make $\eta_0$ small enough for its
reserved norm slack as well. This order is possible because all counts
and density constants are already fixed and independent of these late
grids and margins.

Choose a positive rational $\beta\leq\beta_0$ with
$\beta B_1/3\leq1$ and small enough that
$4\beta B_1/\kappa<d_{\rm sync}/2$. Finally choose positive rational
$\eps<1/10$ and $\delta<1$ so small that
\begin{gather}
 \eps\leq\eps_g,\qquad\eps<\omega/1000,
 \qquad C\delta/\delta_g<q_0/20,\label{out:epsilon-delta}\\
 \frac{4\beta B_1+2C_{\rm ar}\eps
             +6\sqrt{C_{\rm tot}\delta}}{\kappa}<d_{\rm sync},
 \label{out:coarse-budget}\\
 2C_{\rm ar}\eps+6\sqrt{C_{\rm tot}\delta}
                <\beta c_*\kappa/20.
 \label{out:fine-budget}
\end{gather}
There is positive slack in each inequality at $\eps=\delta=0$, so
fixed positive rational choices exist. If necessary decrease them for
any remaining direct numerical slack. All unstable events other than
failed direct cones, including the mapped extraction errors, now total
less than $q_0$.

The finitely many lower-size cutoffs are imposed after these choices by
increasing $C_{\rm base}$. These include $L_{\rm port}$, the finitely
many $L_B(a)$ for the chosen tester accuracies, and $L_g$ for the fixed
arity, size factor, and output loss. This does not create arity feedback. The raw
algorithm's exponent is independent of correction rounds; the chosen
port allowance has slack above that exponent and the degree-ten output
polynomial evaluation. For the fixed cones, increasing the additive
padding absorbs remaining multiplicative factors and cutoffs into a
single fixed size factor, even factor one if desired. Padding checks
and zero construction are $O(L)$ scans with uniform coefficients; their
implementation does not perform hidden work growing with
$C_{\rm base}$. Public computation of bit-polynomial coefficients may
have a larger fixed polynomial exponent than the circuit evaluating
the resulting monomials. Only the latter exponent is charged to the
functional size allowance. No parameter varies with $N$ after this
single choice.

\begin{lemma}[Coarse synchronization]\label{out:coarse-sync}
Every assignment satisfying all but $\delta|T|$ gates has
$\norm{X-\bar X}_{\rm rep}<d_{\rm sync}$, without any assumption that
its functional decisions are correct.
\end{lemma}
\begin{proof}
Apply \eqref{out:coarse-pointwise} to nonfailed direct cones. At failed
ones, both $X_{ir}$ and $(A_{\rm rep}X)_{ir}$ belong to $[-1,2]$, so
their difference is at most three. The factor four in the norm and
Lemma~\ref{out:charges}(i) give
\[
 \norm{X-A_{\rm rep}X}_{\rm rep}
 \leq4\beta B_1+2C_{\rm ar}\eps
                  +6\sqrt{C_{\rm tot}\delta}.
\]
On the other hand \eqref{out:replica-gap} and the reverse triangle
inequality give
$\kappa\norm{X-\bar X}_{\rm rep}
\leq\norm{X-A_{\rm rep}X}_{\rm rep}$.
Equation~\eqref{out:coarse-budget} proves the assertion. In particular,
failures at physical feedback outputs, threshold constants, or the
numeric averaging gates are already included; no decision is trusted
in establishing this estimate.
\end{proof}

\subsection{The clipped mean contradiction and universal decoding}

\begin{lemma}[Fine residual]\label{out:fine-residual}
If an assignment meets the gate-failure budget and
$\dist(\bar X,\mathcal E)>10\tau$, then
\begin{equation}\label{out:residual}
 E_{\rm res}:=
 \norm{X-\clip_{\mathcal B}(A_{\rm rep}X+\beta v(\bar X))}_{\rm rep}
                  <\beta c_*\kappa/2.
\end{equation}
Here clipping is componentwise on the replicas, and $v(\bar X)$ is
repeated identically over them.
\end{lemma}
\begin{proof}
Fix the actual mean $\bar X$ in all ideal tape experiments. Coarse
synchronization permits all spoilage charges already budgeted. Discard
cohorts with a high-failure or high-spoil invocation. In every remaining
invocation whose ideal branch-pattern inputs are all near ports or
public substitutions, Theorem~\ref{thm:port-gadget} applies. Mark all
cohort outputs receiving an incorrect designated wire from any such
invocation. Equations~\eqref{out:designated-count} and
\eqref{out:round-robin} bound their union by
$2B_{\rm out}\upsilon$ of the outputs. Do this marking for every
near-input hypothetical invocation before following the selected path.
Non-near-input invocations are irrelevant unless selected; on the ideal
good-port event a selected invocation necessarily is a near-input one.

Outside these marks and the other unstable events, every direct
comparison in the complete tree has the needed margin and every
selected invocation supplies its correct near-bit. Induct on the
ideal answer sequence: earlier reliable bits select its actual
substitution/function branch, the selected gadget is reliable, and
Boolean shielding removes arbitrary unused outputs. Thus the whole
selected direct computation equals its ideal Boolean run, and
\eqref{out:fine-pointwise} holds.

For an unstable but entirely satisfied direct cone, use only
\eqref{out:coarse-pointwise} and $|v_i(\bar X)|\leq B_1$:
its discrepancy from the clipped true-field update is at most
$3\beta B_1+C_{\rm ar}\eps$. This is proportional to $\beta$ apart
from numeric error. Failed direct cones use only the range bound three.
The empirical ideal error, clipping nonexpansiveness, and the total
unstable fraction at most $q_0$ therefore give
\[
 E_{\rm res}
 \leq\beta\bigl(c_*\kappa/20+6B_1\sqrt{q_0}\bigr)
      +2C_{\rm ar}\eps+6\sqrt{C_{\rm tot}\delta}
 <\beta c_*\kappa/2
\]
by \eqref{out:q0} and \eqref{out:fine-budget}. This includes arbitrary
concentration on extraction gates: above-budget invocations lose their
whole cohort; below-budget invocations lose their guaranteed fraction
of designated wires, and the lower bound in
\eqref{out:designated-count} prevents a concentration factor.
\end{proof}

\begin{lemma}[Clipping and replica means]\label{out:clip-mean}
Let $A$ be a stochastic replica operator satisfying
\eqref{out:replica-gap}, and let $c$ be constant across the replicas of
each component. Put $F=\clip_{\mathcal B}(AX+c)$ and
$E=\norm{X-F}_{\rm rep}$. Then
\[
 \norm{X-JX}_{\rm rep}\leq E/\kappa,\qquad
 \norm{JX-\clip_{\mathcal B}(JX+c)}_{\rm b}\leq E/\kappa.
\]
\end{lemma}
\begin{proof}
Write $Q=I-J$ and $g=\clip_{\mathcal B}(JX+c)$, a repeated vector.
The mean projection minimizes distance to repeated vectors, and
coordinate clipping is nonexpansive. Hence
\[
 \norm{QF}_{\rm rep}
 \leq\norm{F-g}_{\rm rep}
 \leq\norm{A(X-JX)}_{\rm rep}
 \leq(1-\kappa)\norm{QX}_{\rm rep}.
\]
Projecting $X-F$ gives
$\norm{QX}_{\rm rep}\leq E+(1-\kappa)\norm{QX}_{\rm rep}$,
which proves the first inequality. Taking means and then comparing to
$g$ gives
\[
 \norm{JX-g}_{\rm b}
 \leq\norm{J(X-F)}_{\rm rep}+\norm{JF-g}_{\rm rep}
 \leq E+(1-\kappa)\norm{QX}_{\rm rep}\leq E/\kappa.
\]
This argument does not exchange averaging and clipping.
\end{proof}

Apply Lemma~\ref{out:clip-mean} with $c=\beta v(\bar X)$ and the fine
residual. If the mean were farther than $10\tau$ from $\mathcal E$, it
would give
\[
 \norm{\bar X-\clip_{\mathcal B}(\bar X+\beta v(\bar X))}_{\rm b}
 <\beta c_* ,
\]
contrary to \eqref{out:geometry-displacement}. Thus every assignment
meeting the gate-failure budget satisfies
\begin{equation}\label{out:endpoint-proximity}
 \dist(\bar X,\mathcal E)\leq10\tau.
\end{equation}
Choose an endpoint witnessing this bound. Its first block is exactly
$E_w$ for a desired endpoint. Round the first block of $\bar X$ at
$1/2$. Each erroneous bit contributes at least $1/4$ to its squared
coordinate distance from $E_w$, and so the resulting word $y$ obeys
\[
 \dist(y,E_w)\leq4(10\tau)^2<\rho_{\rm dec}.
\]
The global decoder following \eqref{out:binary-encoding} recovers $w$.
The endpoint decoder in Theorem~\ref{thm:outer-representation} and
Lemma~\ref{lem:macrograph} returns a nonzero End-of-Line solution. It
does not require knowing which gates failed or which endpoint witnessed
\eqref{out:endpoint-proximity}.

This defines a deterministic decoder on every acceptable rational
assignment. Means and comparisons can be computed exactly: forming a
common denominator within a group increases its bit length by at most
the sum of its input denominator bit lengths, followed by polynomial
integer-arithmetic overhead. The port and raw algorithms are polynomial
time, with arbitrary fixed defaults on out-of-promise inputs. The
proved proximity places this actual input inside their guaranteed
radii. Hence decoding is polynomial in $N$ plus the total assignment
encoding length. The soundness argument itself applied to all real
assignments satisfying the failure budget.

\subsection{Complete output length, uniformity, and totality}

Every constituent scale is a fixed polynomial in $L$, and
$|\Gamma|=O(N\ell_N^{O(1)})$. Equations~\eqref{out:gate-count} and
\eqref{out:cohort-size} give, for some fixed exponent $b_1$,
\begin{equation}\label{out:size}
 |V|+|T|=O(mT_{\rm tap}K_c)
       =O(|\Gamma|S_0T_{\rm tap}G_0)
       =O(N\ell_N^{b_1}).
\end{equation}
This count includes every copied direct cone, every entire shared
invocation, its extraction gates, every Constant gate, and every
comparison in every hypothetical branch. Routing creates one connection
record per counted input occurrence; it is not an implicit description
of a larger circuit. Explicit node and gate identifiers require
$O(\log(N\ell_N^{b_1}))=O(\ell_N)$ bits. Field/index metadata, locally
generated function coefficients, and every rational parameter have
fixed polynomial-in-$L$ bit bounds. The offsets, bin widths,
$\beta,\eps,\delta$ and numerical precisions are fixed rational
constants. In particular, the threshold-input length qualification in
Theorem~\ref{thm:port-gadget} contributes only a fixed bound here.
If a representation repeats local parameter descriptions instead of
sharing them, multiplying \eqref{out:size} by their polynomial-in-$L$
length still changes only a fixed logarithmic exponent. Input/output
lists and delimiters obey the same bound. Thus there are a fixed
positive integer $A$ and fixed nonnegative integer $b$ for which the
complete output encoding has length at most $AN\ell_N^b$.

All tape masses and their integer multiplicities are deterministically
enumerable. So are every raw-neighborhood and correction-cone address,
every padded function, every inner invocation gate, every replica move,
and every round-robin routing record. Public outer encodings and the
source data are constructed by the polynomial-time algorithms already
proved. Polynomially many operations on polynomial-bit integers suffice
for all these tasks. Products and compositions of the finitely many
fixed polynomial bounds give one deterministic polynomial-time
reduction. It never lists the graph of all logical vertex labels. The
gate set is nonempty, all parameters lie in $[0,1]$, fan-in is at most
two, and the fresh-output construction enforces the required unique
output ownership despite the feedback cycles.

It remains to establish the finite totality assertion for arbitrary
generalized circuits, independently of this reduction.

\begin{lemma}[A fixed rational mesh satisfies every gate]
\label{out:totality}
For every fixed rational $0<\eps<1/10$, every generalized circuit in
Definition~\ref{def:gcircuit} has a rational assignment satisfying all
its gates, each coordinate having denominator
$q=\lceil20/\eps\rceil$. The complete assignment length is bounded by
one fixed polynomial with nonnegative integer coefficients in the
circuit encoding length.
\end{lemma}
\begin{proof}
Construct a continuous self-map of $[0,1]^V$. At an arithmetic gate's
unique output use its exact numerical operation. At a Less gate with
inputs $a,b$, use the clipped linear ramp that is zero when
$x_b-x_a\leq-\eps/2$ and one when
$x_b-x_a\geq\eps/2$. For Boolean gates define a continuous signal
$h:[0,1]\to[0,1]$ equal to zero on $[0,2\eps]$, equal to one on
$[1-2\eps,1]$, and linear between them. Use respectively
\[
 \max\{h(x_a),h(x_b)\},\qquad
 \min\{h(x_a),h(x_b)\},\qquad 1-h(x_a)
\]
for Or, And, and Not. Set coordinates not produced by gates to zero.
Unique output ownership makes these prescriptions a well-defined
continuous cube map. Brouwer's fixed-point theorem~\cite{Brouwer1911} gives a fixed point
$x^*$; no efficient computation of it is asserted.

Round every coordinate to the nearest point of
$\{0,1/q,\ldots,1\}$, obtaining $\widetilde x$ with
$|\widetilde x_v-x_v^*|\leq1/(2q)\leq\eps/40$.
Constant errors are at most $\eps/40$, Scale and Copy errors at most
$2\eps/40$, and Add and Subtract errors at most $3\eps/40$, using
parameter bounds and nonexpansive clipping. For a triggered strict
Less antecedent after rounding, its signed input gap before rounding
has the same sign and magnitude strictly greater than
$\eps-2\eps/40>\eps/2$. Its fixed-point output is therefore the
required endpoint zero or one, which rounding preserves.

If a weak Boolean low antecedent
$\widetilde x_a\leq\eps$ holds, then
$x_a^*\leq\eps+\eps/40<2\eps$, so $h(x_a^*)=0$.
Similarly $\widetilde x_a\geq1-\eps$ implies
$h(x_a^*)=1$. Every applicable Or, And, or Not conclusion follows
from its displayed exact formula, including equality at the weak
threshold. The corresponding fixed-point output is zero or one and
is unchanged by rounding. This verifies every gate implication.

The denominator $q$ is a fixed integer depending only on $\eps$.
There are at most a fixed multiple of $M+1$ explicitly listed nodes
in a length-$M$ input, and each output rational has a fixed bit bound.
Accounting also for order, indices if used, and delimiters, a
sufficiently large fixed integer $C$ makes
$p(M)=C(M+1)^2$ an upper bound for the entire assignment. This
polynomial has nonnegative integer coefficients and is independent of
the particular circuit and its parameter denominators.
\end{proof}

\begin{proof}[Proof of Theorem~\ref{thm:main}]
The preceding argument uses Theorem~\ref{thm:port-gadget}, whose construction
is proved in Sections~\ref{sec:inner-algebra}--\ref{sec:inner-gadget}, and
Lemma~\ref{lem:geometric-procedures}, proved in Section~\ref{sec:geometry}.
Use the fixed rational $\eps,\delta$ chosen in
\eqref{out:epsilon-delta}--\eqref{out:fine-budget} and the polynomial
$p$ from Lemma~\ref{out:totality}. That lemma proves totality, with all
gates satisfied. The explicit construction and \eqref{out:size} give
the claimed complete quasilinear encoding bound and a deterministic
polynomial-time reduction. Equations~\eqref{out:endpoint-proximity}
and~\eqref{out:tau}, followed by the global endpoint decoder, give the
required solution from every acceptable assignment, with the stated
polynomial running time. Every constant and algorithm has been fixed
independently of the input. This proves all four assertions.
\end{proof}
